\pagestyle{mystyle}
\chapter{Ottimizzazione e Adattamento al Dominio}
\label{chap:ottimizzazione}

Il capitolo analizza le strategie adottate per adattare il sistema RAG al
dominio i3Guild senza trasformare il prototipo in un progetto di addestramento
del modello. Questa scelta deriva da un vincolo concreto: non è disponibile un
glossario tecnico aziendale validato e non esiste, allo stato attuale, un
dataset supervisionato sufficiente per un fine-tuning robusto. Il corpus
utilizzabile è costituito principalmente dai dati delle Gilde, dai relativi
metadati, dalle descrizioni testuali e dai contenuti pubblicati nel tempo.

L'adattamento al dominio viene quindi affrontato in modo non parametrico. Il
sistema non modifica i pesi del modello, ma interviene sulla qualità del
retrieval, sulla normalizzazione delle query, sull'uso dei metadati, sulla
costruzione del prompt e sulla configurazione dei modelli locali. In questo
modo il sistema rimane aggiornabile, eseguibile on-premise e coerente con i
vincoli architetturali definiti nei capitoli precedenti.

La prima parte del capitolo descrive i vincoli sperimentali: assenza di un
glossario formale, disponibilità parziale dei dati aziendali, risorse hardware
limitate e necessità di mantenere separati il sistema applicativo e la
componente RAG. La seconda parte approfondisce le ottimizzazioni implementate:
query rewriting, stopword di dominio, filtri strutturati, deduplicazione,
chunking, soglie dinamiche di retrieval e gestione del contesto generativo. La
terza parte riguarda l'estensione agentica, trattata come evoluzione del flusso
RAG tradizionale e non come sostituzione dell'architettura di base.

Il capitolo include inoltre una sezione di sperimentazione preliminare su
strategie alternative di adattamento. Tecniche come LoRA, QLoRA,
fine-tuning e interventi a tempo di inferenza vengono considerate in termini di
fattibilità, requisiti e limiti, senza renderle parte del sistema aziendale
integrato. Questa analisi serve a delimitare il perimetro del prototipo e a
motivare perché, nel lavoro svolto, l'adattamento non parametrico rappresenti la
scelta più solida.

La chiusura del capitolo definisce la configurazione finale da portare alla
valutazione del Capitolo~\ref{chap:evaluation}: parametri di retrieval,
criteri di selezione dei documenti, impostazioni del modello locale e modalità
agentica. In questo modo il Capitolo 5 può concentrarsi sulla valutazione del
sistema risultante, evitando di riaprire decisioni architetturali già motivate.

% Struttura prevista del capitolo:
% 1. Vincoli e perimetro dell'ottimizzazione
% 2. Adattamento non parametrico al dominio i3Guild
% 3. Ottimizzazione del retrieval
% 4. Ottimizzazione del prompt e della generazione locale
% 5. Estensione agentica come adattamento del flusso utente
% 6. Sperimentazione preliminare su tecniche alternative
% 7. Configurazione finale per la valutazione
